\section{자를 올리는 조건에 크기를 넣는다}

\claimx{앞 노트가 정본 자로 올린 도메인 균등 가중은 학습 라벨을 전량 부수어도 자기 표준오차의 두 배를 못 넘는다 --- 그래서 승격을 되돌린다}{그 자의 라벨 전량 파괴 델타가 등록된 정칙화 두 값 모두에서 자기 표준오차의 두 배를 넘으면 이 주장은 떨어진다.}

앞 노트의 채택 규칙은 \textbf{부호만} 물었다. 라벨을 전부 부수면 자가 내려가는가 --- 부호가 음수인가 --- 만 보고 자를 올렸다. 이 노트는 거기에 \textbf{크기}를 넣는다: 그 델타가 자기 표준오차의 두 배를 넘는가, 그리고 등록된 문턱을 넘는가. 두 정칙화 값 모두에서 물어야 한다.

결과는 등록된 정칙화 두 값에서 각각 $1.9527$ 배와 $0.3304$ 배다. 여덟 칸 파괴 대조 전체로 넓히면 균등 가중은 $0 / 8$ 이고 앞 노트가 곁으로 내린 행 가중은 $3 / 8$ 다.

\section{까닭 --- 균등평균이 잡음을 증폭한다}

\claimx{단순 균등평균은 표본이 작은 도메인과 큰 도메인에 같은 무게를 주어 자의 표준오차를 넓힌다 --- 그것이 검정력이 영이 되는 기전이다}{도메인 안 순열 귀무의 표준편차가 도메인 크기와 무관하면 이 설명은 떨어진다.}

도메인마다 유보 정답을 그 도메인 안에서만 뒤섞어 순위상관의 귀무 표준편차를 냈다. 이 값은 그 도메인의 행 수와 동률 구조만의 함수이고 예측을 보지 않는다. 가장 작은 도메인과 가장 큰 도메인의 귀무 표준편차 비는 $8.1219$ 배다. 균등평균은 그 둘을 같은 무게로 놓는다.

그래서 자 넷을 한 족으로 놓고 잰다. 행 가중, 균등, 귀무 표준편차의 역수 가중, 귀무 분산의 역수 가중이다. 가장 큰 도메인이 차지하는 무게는 각각 $0.545994$, $0.1$, $0.278683$, $0.542072$ 다.

\textbf{여기서 자기 적발이 하나 나온다.} 분산 역수 가중은 최소분산 가중이지만 귀무 분산이 행 수의 역수에 비례하므로 \textbf{행 가중으로 되돌아간다} --- 지배 도메인을 덜어내려던 처방이 원래 자로 수렴한다. 그 사이에 서는 것이 \textbf{표준편차의 역수 가중}이다.

\section{같은 조건에 두 값이 커밋된 까닭}

\claimx{앞 노트의 산출물 두 개가 같은 이름의 조건에 서로 다른 부호를 적은 것은 수가 틀려서가 아니라 등록된 조건이 뽑기 수를 안 적었기 때문이다}{두 값 중 하나라도 자기 설정에서 재현되지 않으면 이 설명은 떨어진다.}

한쪽은 한 벌에서 $0.086077$, 다른 쪽은 스물다섯 벌 평균에서 $-0.023854$ 였다. 각자의 설정에서 다시 내니 $0.086077$ 와 $-0.023854$ 로 \textbf{둘 다 소수 여섯 자리까지 재현된다}. 한 벌 추정의 분포를 이백 벌에서 내니 표준편차가 $0.078638$ 이고 부호가 양수인 비율이 $0.38$ 다. \textbf{조건의 답을 정한 것은 자료가 아니라 뽑기다.}

\section{크기를 맞춘 대조}

\claimx{안쪽 원천의 학습 라벨을 부수는 것이 바깥 원천에서 같은 수의 행을 부수는 것보다 자를 더 움직인다 --- 비가 아니라 크기 맞춤 대조로 말한다}{같은 행 수 대조의 두 절대 델타 차가 짝 표준오차의 두 배를 못 넘으면 이 주장은 떨어진다.}

앞 노트는 행당 배율을 헤드라인에 실었는데 그 분모는 바깥 원천 라벨 전량 파괴의 델타였고, 그 값은 영을 무는 잡음이며 \textbf{부호가 양수}였다 --- 부수면 좋아진다는 뜻이다. 비를 버린다. 같은 행 수를 부수는 두 대조만 쓴다. 안쪽 대조는 $-0.113589 \pm 0.052258$ 이고 바깥 대조는 $0.000473 \pm 0.009252$ 이며 그 차는 짝 표준오차의 $2.4969$ 배다. 다른 정칙화 값에서는 안쪽 대조 자체가 $1.4891$ 배로 \textbf{등록 문턱에 못 미친다} --- 그 사실도 같이 적는다.

행당 환산을 그래도 살려 쓸 자리가 하나 있어 보였다. 바깥 원천 안에서 행 수만 다른 두 대조의 비가 행 수 비에 가깝다는 것이다. 앞 노트의 산출물에서 읽으면 그 비는 $20.6397$ 이고 행 수 비는 $19$ 다. 	extbf{그런데 벌 수 규약을 그대로 두고 뒤섞기 씨앗의 밑값만 바꾸면} 그 비가 $52.6533$ 로 $2.5511$ 배 움직인다. 분모가 영 근처의 잡음이기 때문이다 --- 	extbf{행당 환산이 거의 선형이라는 관찰은 한 씨앗에서만 참이다.}

\section{특징 쪽을 부순다}

\claimx{바깥 원천의 특징 행렬을 통째로 뒤섞어도 자가 잡음만큼도 안 움직인다 --- 라벨 쪽 결과와 같은 꼴이다}{그 델타가 자기 표준오차의 두 배를 넘으면 이 주장은 떨어진다.}

앞 노트는 라벨만 부쉈다. 이 노트는 학습 행의 특징 행렬을 같은 구간에서 뒤섞는다. 바깥 원천 특징 전량을 뒤섞은 델타는 $0.024195 \pm 0.064171$ 이고 안쪽 원천 쪽은 $-0.108357$ 다. 그리고 바깥 원천 행은 게이트 도메인 $10 / 10$ 을 덮는다 --- \textbf{도메인이 안 겹쳐서 안 쓰인다는 설명은 안 선다.}

